%% toc: P15 161 161.1
% Proofs for mean-power convergence from convergence in probability, and for the source forms of the measure foundations.
\section{Mean-power convergence from convergence in probability: proofs}
\label{sec:annex-p15-w02c-p15-014-probability-mean}

Use the definitions of Chapter~\ref{sec:p15-w02c-p15-014-probability-mean}
and the notation \(V_n=|X_n-X|^p\), \(W_n=|X_n|^p\), \(W=|X|^p\) of
\cref{def:p15-w02c-p15-014-reading}. These are nonnegative, finite and
measurable, as continuous functions of finite real random variables
(\cref{thm:p15-measurable-operations}). For nonnegative measurable \(U\),
expectations are additive over a partition of \(\Omega\) into finitely
many events and are monotone in \(U\)
(\cref{thm:p15-nonnegative-integral-laws}~(2) and~(4)).

\begin{annexitem}{annex:p15-w02c-p15-014-prob-mean}{thm:p15-w02c-p15-014-probability-mean}{Tails and moment hypotheses}
\begin{proof}
\emph{Step 0: four facts.}

(a) For \(\eta>0\), \(\{V_n>\eta\}=\{|X_n-X|>\eta^{1/p}\}\), because
\(t\mapsto t^p\) is strictly increasing on \([0,\infty)\). Hence
\(\mathbb P(V_n>\eta)\to0\) for every \(\eta>0\) whenever
\(X_n\to X\) in probability.

(b) By \cref{thm:p15-w02b-p15-013-power}~(3), applied pointwise to the
pairs \((X_n,-X)\), \((X_n-X,X)\) and \((X_n,X-X_n)\),
\[
 V_n\leq c_p(W_n+W),\qquad W_n\leq c_p(V_n+W),\qquad
 W\leq c_p(W_n+V_n),
\]
where \(c_p=\max(1,2^{p-1})\).

(c) Let \(J\) be an index set and let \((A_j)_{j\in J}\) and
\((B_j)_{j\in J}\) be uniformly integrable families of nonnegative
variables. If
\(0\leq U_j\leq c(A_j+B_j)\) with measurable \(U_j\) and a constant
\(c>0\), then \((U_j)_{j\in J}\) is uniformly integrable. Indeed, each
\(U_j\) is integrable by monotonicity, and by
\cref{thm:p15-ui-small-sets} it suffices to check bounded means and
small-event control. First,
\(\sup_j\mathbb EU_j\leq c(\sup_j\mathbb EA_j+\sup_j\mathbb EB_j)<\infty\).
Second, given \(\epsilon>0\), small-event control of the two families
gives \(\delta>0\) such that \(\mathbb P(A)\leq\delta\) implies
\(\mathbb E(A_j;A)\leq\epsilon/(2c)\) and
\(\mathbb E(B_j;A)\leq\epsilon/(2c)\) for every \(j\); then
\(\mathbb E(U_j;A)\leq c\,\mathbb E(A_j;A)+c\,\mathbb E(B_j;A)\leq\epsilon\).
A family each of whose members is one fixed integrable variable is
uniformly integrable, being dominated by that variable
(\cref{thm:p15-ui-small-sets}).

(d) If \((U_j)_{j\in J_1}\) and \((U_j)_{j\in J_2}\) are uniformly
integrable, so is \((U_j)_{j\in J_1\cup J_2}\): for every \(M\),
\(\sup_{j\in J_1\cup J_2}\mathbb E(|U_j|;|U_j|>M)\) is the larger of
the two suprema over \(J_1\) and \(J_2\), and both tend to zero as
\(M\to\infty\). Every finite family of integrable variables is
uniformly integrable (\cref{thm:p15-ui-small-sets}).

\emph{Item (1), necessity.} Suppose \(\mathbb EV_n\to0\) in the sense
of \cref{def:p15-w02c-p15-014-reading}; convergence in probability is
not used in this part. Taking \(\varepsilon=1\) there gives \(N\) with
\(\mathbb EV_n\leq1<\infty\) for all \(n\geq N\). We show that
\((V_n)_{n\geq N}\) is uniformly integrable. Let \(\epsilon>0\). Choose
\(N'\geq N\) with \(\mathbb EV_n\leq\epsilon\) for every \(n>N'\). The
finite family \(V_N,\ldots,V_{N'}\) of integrable variables is uniformly
integrable, so there is \(M_0\) with
\(\mathbb E(V_n;V_n>M)\leq\epsilon\) for \(N\leq n\leq N'\) and
\(M\geq M_0\). For \(n>N'\) and every \(M\), monotonicity gives
\(\mathbb E(V_n;V_n>M)\leq\mathbb EV_n\leq\epsilon\). Hence
\(\sup_{n\geq N}\mathbb E(V_n;V_n>M)\leq\epsilon\) for all \(M\geq M_0\).
As \(\epsilon\) was arbitrary, the supremum tends to zero as
\(M\to\infty\). If moreover \(\mathbb EV_n<\infty\) for every \(n\),
the finite family \(V_1,\ldots,V_{N-1}\) is uniformly integrable, and
by fact~(d) so is \((V_n)_{n\geq1}\).

\emph{Item (1), sufficiency.} Suppose \(X_n\to X\) in probability and
that \((V_n)_{n\geq N}\) consists of integrable variables and is
uniformly integrable. Let \(\epsilon>0\) and choose \(M>0\) with
\(\sup_{n\geq N}\mathbb E(V_n;V_n>M)\leq\epsilon\). Put
\(\eta=\min(\epsilon,M/2)\), so \(0<\eta<M\). For \(n\geq N\), the
events \(\{V_n\leq\eta\}\), \(\{\eta<V_n\leq M\}\) and \(\{V_n>M\}\)
partition \(\Omega\). On the first \(V_n\leq\eta\), and on the second
\(V_n\leq M\). Therefore
\[
 \begin{aligned}
 \mathbb EV_n&\leq\eta\,\mathbb P(V_n\leq\eta)
     +M\,\mathbb P(\eta<V_n\leq M)+\mathbb E(V_n;V_n>M)\\
 &\leq\epsilon+M\,\mathbb P(V_n>\eta)+\epsilon .
 \end{aligned}
\]
By fact~(a), \(\mathbb P(V_n>\eta)\to0\), so there is \(N''\geq N\) with
\(M\,\mathbb P(V_n>\eta)\leq\epsilon\) for \(n\geq N''\), and then
\(0\leq\mathbb EV_n\leq3\epsilon\). As \(\epsilon\) was arbitrary,
\(\mathbb EV_n\to0\), that is, \(X_n\to X\) in \(L^p\). If
\(\mathbb EV_n<\infty\) for every \(n\), the condition with \(N=1\) is
the second sentence of item~(1).

\emph{Item (2).} Assume \(\mathbb EW_n<\infty\) for every \(n\).

Suppose first that \((W_n)_{n\geq1}\) is uniformly integrable. The map
\(x\mapsto|x|^p\) is continuous on \(\mathbb R\), so
\(W_n\to W\) in probability by
\cref{thm:p15-w02b-p15-013-as-probability}~(4). By
\cref{thm:p15-probability-l1-tools}~(4) and
\cref{thm:p15-ui-small-sets},
\(\mathbb EW\leq\liminf_n\mathbb EW_n\leq\sup_n\mathbb EW_n<\infty\);
so \(X\in L^p\). The constant family \((W)_{n\geq1}\) is uniformly
integrable by fact~(c), and the first inequality of fact~(b) with
fact~(c) shows that \((V_n)_{n\geq1}\) consists of integrable variables
and is uniformly integrable. The sufficiency half of item~(1), with
\(N=1\), gives \(X_n\to X\) in \(L^p\).

Suppose conversely that \(X_n\to X\) in \(L^p\). By the necessity half
of item~(1) there is \(N\) such that \(V_n\) is integrable for
\(n\geq N\) and \((V_n)_{n\geq N}\) is uniformly integrable. The third
inequality of fact~(b), at \(n=N\), gives
\(\mathbb EW\leq c_p(\mathbb EW_N+\mathbb EV_N)<\infty\), so
\(X\in L^p\). For \(n\geq N\) the second inequality of fact~(b) and
fact~(c), applied to the families \((V_n)_{n\geq N}\) and the constant
family \(W\), show that \((W_n)_{n\geq N}\) is uniformly integrable. The
finite family \(W_1,\ldots,W_{N-1}\) is integrable by hypothesis, hence
uniformly integrable, and fact~(d) shows that \((W_n)_{n\geq1}\) is
uniformly integrable. In both directions \(X\in L^p\).

\emph{Item (3).} Assume \(\mathbb EW<\infty\).

If \(X_n\to X\) in \(L^p\), the necessity half of item~(1) gives \(N\)
such that \((V_n)_{n\geq N}\) consists of integrable variables and is
uniformly integrable. For \(n\geq N\), \(W_n\leq c_p(V_n+W)\) by
fact~(b), so \(W_n\) is integrable, that is \(X_n\in L^p\), and
\((W_n)_{n\geq N}\) is uniformly integrable by fact~(c).

Conversely, let \(N\) be such that \(X_n\in L^p\) for \(n\geq N\) and
\((W_n)_{n\geq N}\) is uniformly integrable. Put
\(X'_k=X_{N+k-1}\) for \(k\geq1\). Then \(X'_k\to X\) in probability,
since for each \(\varepsilon>0\) the sequence
\(\mathbb P(|X'_k-X|>\varepsilon)\) is a tail of a sequence tending to
zero; every \(X'_k\) lies in \(L^p\); and \((|X'_k|^p)_{k\geq1}\) is
the family \((W_n)_{n\geq N}\). Item~(2) gives
\(\mathbb E|X'_k-X|^p\to0\) as \(k\to\infty\), and therefore
\(\mathbb E|X_n-X|^p\to0\) as \(n\to\infty\), because a sequence
converges exactly when one of its tails converges, to the same limit.

\emph{Item (4).} Sufficiency of the condition is the first half of the
proof of item~(2). For the variable \(Y_p\), let \(Y\) be the variable
of \cref{thm:p15-w02b-p15-013-countermodels}~(5): on the uniform space,
\(Y=2^k\) on \([2^{-k},2^{1-k})\) for \(k\geq1\). Each point of
\((0,1)\) lies in exactly one of these intervals, so \(Y\) is finite and
positive, and \(\mathbb EY=\infty\) is proved there. Put
\(Y_p\defeq Y^{1/p}\). It is finite, positive and measurable, as a
continuous function of \(Y\), and
\(\mathbb EY_p^{\,p}=\mathbb EY=\infty\).

(a) Let \(X_n=X=Y_p\). Then \(|X_n-X|=0\) everywhere, so
\(\mathbb P(|X_n-X|>\varepsilon)=0\) for every \(\varepsilon>0\) and
\(\mathbb E|X_n-X|^p=0\) for every \(n\): the sequence converges to
\(X\) in probability and in \(L^p\). But \(\mathbb E|X_n|^p=\infty\)
for every \(n\), so no \(X_n\) lies in \(L^p\), and
\((|X_n|^p)_{n\geq1}\) is not a family of integrable variables, hence
not uniformly integrable (\cref{def:p15-w02c-p15-014-reading}). Here
\(X\notin L^p\), and item~(1) holds with \(V_n=0\).

(b) Let \(X=0\), \(X_1=Y_p\) and \(X_n=0\) for \(n\geq2\). For
\(\varepsilon>0\) and \(n\geq2\), \(\mathbb P(|X_n-X|>\varepsilon)=0\),
and \(\mathbb E|X_n-X|^p=0\); so \(X_n\to X\) in probability and in
\(L^p\), although \(\mathbb E|X_1-X|^p=\mathbb EY=\infty\). Here
\(X\in L^p\), \(X_1\notin L^p\), so \((|X_n|^p)_{n\geq1}\) is not
uniformly integrable, while the tail from \(N=2\) consists of the zero
variable and is uniformly integrable, as item~(3) requires.
\end{proof}
\end{annexitem}

\section{Source forms in the measure foundations: proofs}
\label{sec:annex-p15-w02c-p15-014-foundations}

\begin{annexitem}{annex:p15-w02c-p15-014-field}{thm:p15-w02c-p15-014-sigma-field-variant}{Nonemptiness suffices}
\begin{proof}
(1) Suppose \(\mathcal F\) is a sigma-algebra in the sense of
\cref{def:p15-sigma-algebra}. It contains \(\Omega\), so it is nonempty,
and it is closed under complements and countable unions by definition.

Conversely, suppose \(\mathcal F\) is nonempty and satisfies~(a)
and~(b). Choose \(A\in\mathcal F\). By~(a), \(A^c\in\mathcal F\). The
two-term sequence \(A,A^c\) is countable in the stated sense, so by~(b)
its union \(A\cup A^c\) lies in \(\mathcal F\). Every outcome lies in
\(A\) or in \(A^c\), so \(A\cup A^c=\Omega\), and \(\Omega\in\mathcal F\).
Together with~(a) and~(b), this is \cref{def:p15-sigma-algebra}.

(2) Suppose \(\mathcal A\) is an algebra in the sense of
\cref{def:p15-sigma-algebra}: it contains \(\Omega\) and is closed
under complements and finite unions. Then it is nonempty, and for
\(A,B\in\mathcal A\) the union \(A\cup B\) is a finite union of members,
so \(A^c\in\mathcal A\) and \(A\cup B\in\mathcal A\).

Conversely, let \(\mathcal A\) be nonempty with the two closure
conditions. Choose \(A\in\mathcal A\); then \(A^c\in\mathcal A\) and
\(\Omega=A\cup A^c\in\mathcal A\). Closure under finite unions: the
empty union is \(\varnothing=\Omega^c\in\mathcal A\); a union of one
member is that member; and if every union of \(k\geq1\) members lies in
\(\mathcal A\), then for members \(A_1,\ldots,A_{k+1}\) the union
\(\bigcup_{i\leq k+1}A_i=\bigl(\bigcup_{i\leq k}A_i\bigr)\cup A_{k+1}\)
is a union of two members of \(\mathcal A\) and lies in \(\mathcal A\).
By induction every finite union of members lies in \(\mathcal A\), so
\(\mathcal A\) is an algebra.

For intersections, let \(A,B\in\mathcal A\). Then
\(A^c,B^c\in\mathcal A\), so \(A^c\cup B^c\in\mathcal A\) and its
complement lies in \(\mathcal A\). An outcome lies outside
\(A^c\cup B^c\) exactly when it lies in neither \(A^c\) nor \(B^c\),
that is, in both \(A\) and \(B\); hence
\((A^c\cup B^c)^c=A\cap B\in\mathcal A\).

(3) Each closure condition has the form ``if certain sets belong to the
collection, then certain other sets belong to it''. In the empty
collection no set belongs, so no instance of a premise holds and every
condition is satisfied. The empty collection does not contain \(\Omega\),
so it is neither a sigma-algebra nor an algebra in the sense of
\cref{def:p15-sigma-algebra}.
\end{proof}
\end{annexitem}

\begin{annexitem}{annex:p15-w02c-p15-014-covering}{thm:p15-w02c-p15-014-covering-subadditivity}{Disjointing the cover inside the set}
\begin{proof}
(1) For \(m\geq1\) put \(A'_m=A_m\cap A\), and define
\[
 B_1=A'_1,\qquad
 B_m=A'_m\setminus\bigcup_{k=1}^{m-1}A'_k\quad(m>1).
\]
Each \(B_m\) is obtained from members of \(\mathcal F\) by finitely many
intersections, unions and complements, so \(B_m\in\mathcal F\)
(\cref{thm:p15-generated-events}). By construction
\(B_m\subseteq A'_m\subseteq A_m\) and \(B_m\subseteq A\).

The sets \(B_m\) are pairwise disjoint: if \(k<m\), then
\(B_k\subseteq A'_k\), while \(B_m\) contains no point of \(A'_k\).

Their union is \(A\). Each \(B_m\) lies in \(A\). Conversely let
\(\omega\in A\). Since \(A\subset\bigcup_mA_m\), the set of indices
\(m\) with \(\omega\in A_m\) is nonempty; let \(m_0\) be its least
member. Then \(\omega\in A_{m_0}\cap A=A'_{m_0}\), and for \(k<m_0\) the
point \(\omega\) is not in \(A_k\), hence not in \(A'_k\). So
\(\omega\in B_{m_0}\).

Countable additivity (\cref{def:p15-measure}) applied to the disjoint
sequence \((B_m)\) gives \(\mu(A)=\sum_m\mu(B_m)\). Monotonicity
(\cref{thm:p15-measure-laws}~(1)) gives \(\mu(B_m)\leq\mu(A_m)\) for each
\(m\). Every finite partial sum of the first series is therefore at most
the corresponding partial sum of the second, and the sums, as suprema
of their partial sums in \([0,\infty]\), satisfy
\[
 \mu(A)=\sum_{m=1}^\infty\mu(B_m)\leq\sum_{m=1}^\infty\mu(A_m).
\]

(2) The covering form implies the union form: given any sequence
\((A_m)\) in \(\mathcal F\), the set \(A=\bigcup_mA_m\) lies in
\(\mathcal F\) and is covered by \((A_m)\), so~(1) gives
\(\mu(\bigcup_mA_m)\leq\sum_m\mu(A_m)\). The union form implies the
covering form: if \(A\subset\bigcup_mA_m\), monotonicity gives
\(\mu(A)\leq\mu(\bigcup_mA_m)\), and the union form bounds the right
side by \(\sum_m\mu(A_m)\).

(3) Given \(A\subset A_1\cup\cdots\cup A_k\), put \(A_m=\varnothing\) for
\(m>k\). The sequence covers \(A\), so~(1) applies, and since
\(\mu(\varnothing)=0\) (\cref{def:p15-measure}) the series equals
\(\sum_{m\leq k}\mu(A_m)\).
\end{proof}
\end{annexitem}

\begin{annexitem}{annex:p15-w02c-p15-014-borel-functions}{thm:p15-w02c-p15-014-borel-functions}{Products and composition}
\begin{proof}
(1) Let \(A_1,\ldots,A_n\) be Borel subsets of \(\mathbb R\). An outcome
\(\omega\) satisfies \(V(\omega)\in A_1\times\cdots\times A_n\) exactly
when \(X_i(\omega)\in A_i\) for every \(i\), which gives the displayed
identity. Each \(\{X_i\in A_i\}\) lies in \(\mathcal F\) because \(X_i\)
is a random variable, and a finite intersection of events is an event
(\cref{thm:p15-generated-events}).

Let \(\mathcal P_n\) be the collection of all products
\(A_1\times\cdots\times A_n\) of Borel subsets of \(\mathbb R\). We show
\(\sigma(\mathcal P_n)=\mathcal B(\mathbb R^n)\). For each \(i\) the
coordinate map \(\pi_i(x)=x_i\) from \(\mathbb R^n\) to \(\mathbb R\) is
continuous, hence Borel measurable
(\cref{thm:p15-measurable-operations}), so
\(\pi_i^{-1}(A_i)\in\mathcal B(\mathbb R^n)\), and
\[
 A_1\times\cdots\times A_n=\bigcap_{i=1}^n\pi_i^{-1}(A_i)
 \in\mathcal B(\mathbb R^n).
\]
Thus \(\mathcal P_n\subseteq\mathcal B(\mathbb R^n)\), and since
\(\mathcal B(\mathbb R^n)\) is a sigma-algebra containing
\(\mathcal P_n\), the definition of the generated sigma-algebra gives
\(\sigma(\mathcal P_n)\subseteq\mathcal B(\mathbb R^n)\). Conversely, an
open interval \((a_j,b_j)\) is an open subset of \(\mathbb R\) and hence
a Borel set (\cref{def:p15-borel-sets}), so every open rectangle
\(\prod_{j=1}^n(a_j,b_j)\) with rational \(a_j<b_j\) belongs to
\(\mathcal P_n\). These rectangles generate \(\mathcal B(\mathbb R^n)\)
(\cref{thm:p15-borel-generators}), so
\(\mathcal B(\mathbb R^n)\subseteq\sigma(\mathcal P_n)\).

Since \(V^{-1}(P)\in\mathcal F\) for every \(P\in\mathcal P_n\) and
\(\sigma(\mathcal P_n)=\mathcal B(\mathbb R^n)\), the generator test
(\cref{thm:p15-measurability-criterion}) shows that \(V\) is measurable
from \((\Omega,\mathcal F)\) to \((\mathbb R^n,\mathcal B(\mathbb R^n))\),
that is, a random vector (\cref{def:p15-measurable-map}).

(2) By definition \(f(X_1,\ldots,X_n)(\omega)=f(V(\omega))\), so
\(f(X_1,\ldots,X_n)=f\circ V\). The map \(V\) is measurable by~(1) and
\(f\) is measurable by hypothesis, so their composite is measurable from
\((\Omega,\mathcal F)\) to \((\mathbb R,\mathcal B(\mathbb R))\)
(\cref{thm:p15-measurable-operations}, first sentence). Directly: for
\(B\in\mathcal B(\mathbb R)\),
\(\{f\circ V\in B\}=V^{-1}(f^{-1}(B))\), and
\(f^{-1}(B)\in\mathcal B(\mathbb R^n)\), so this set lies in
\(\mathcal F\). Hence \(f(X_1,\ldots,X_n)\) is a real random variable.

(3) Let \(\Omega=\{1,2\}\), \(\mathcal F=\{\varnothing,\Omega\}\) and
\(\mathbb P(\Omega)=1\); \(\mathcal F\) contains \(\Omega\) and is closed
under complements and unions, so this is a probability space. Let
\(X_1=X_2=\mathbf1_{\{1\}}\), the function equal to one at the outcome 1
and zero at the outcome 2. Then
\(\{X_1\leq1/2\}=\{2\}\notin\mathcal F\), so neither \(X_1\) nor \(X_2\)
is a random variable. The function \(f(x_1,x_2)=x_1-x_2\) is continuous,
hence Borel measurable (\cref{thm:p15-measurable-operations}), and
\(f(X_1,X_2)\) is the constant 0. For every Borel \(B\),
\(\{f(X_1,X_2)\in B\}\) is \(\Omega\) if \(0\in B\) and \(\varnothing\)
otherwise, and both lie in \(\mathcal F\). So \(f(X_1,X_2)\) is a random
variable although no \(X_i\) is.
\end{proof}
\end{annexitem}

\begin{annexitem}{annex:p15-w02c-p15-014-rectangles}{thm:p15-w02c-p15-014-half-open-rectangles}{Intersections, generation and uniqueness}
\begin{proof}
(1) First let \(d=1\) and let \(a,b,a',b'\) be real. A real \(x\) lies
in \((a,b]\cap(a',b']\) exactly when \(a<x\), \(a'<x\), \(x\leq b\)
and \(x\leq b'\), that is, exactly when
\(\max(a,a')<x\leq\min(b,b')\). Hence
\((a,b]\cap(a',b']=(\max(a,a'),\min(b,b')]\), with both sides empty
when \(\max(a,a')\geq\min(b,b')\). For general \(d\), a point
\(x=(x_1,\ldots,x_d)\) lies in \(R(a,b)\cap R(a',b')\) exactly when
\(x_j\in(a_j,b_j]\cap(a'_j,b'_j]\) for every \(j\), that is, by the
case \(d=1\), exactly when \(x\in R(a\vee a',b\wedge b')\). The vectors
\(a\vee a'\) and \(b\wedge b'\) are real, so the intersection belongs
to \(\mathcal H_d\).

(2) \emph{Each half-open rectangle is Borel.} For real vectors \(a,b\)
and \(m\geq1\) let \(G_m=\prod_{j=1}^d(a_j,b_j+1/m)\). Each \(G_m\) is
open: it is empty if some \(b_j+1/m\leq a_j\); otherwise, for
\(x\in G_m\), put
\(r=\min_j\min(x_j-a_j,\,b_j+1/m-x_j)>0\); a point \(y\) with
Euclidean distance \(|y-x|<r\) has \(|y_j-x_j|\leq|y-x|<r\) for every
\(j\), so \(y\in G_m\). Next, \(R(a,b)=\bigcap_{m\geq1}G_m\): if
\(a_j<x_j\leq b_j\) for every \(j\), then \(x_j<b_j+1/m\) for every
\(m\); and if \(a_j<x_j<b_j+1/m\) for every \(m\) and \(j\), then
\(x_j\leq b_j\), since \(x_j>b_j\) would fail for \(1/m<x_j-b_j\).
So \(R(a,b)\) is a countable intersection of open sets, and lies in
\(\mathcal B(\mathbb R^d)\) by \cref{thm:p15-generated-events}. Since
\(\mathcal H_d\subseteq\mathcal B(\mathbb R^d)\), the minimality in
\cref{thm:p15-generated-events} gives
\(\sigma(\mathcal H_d)\subseteq\mathcal B(\mathbb R^d)\).

\emph{Each rational open rectangle lies in \(\sigma(\mathcal H_d)\).}
Let \(Q=\prod_{j=1}^d(a_j,b_j)\) with rational \(a_j<b_j\). We show
\(Q=\bigcup_{m\geq1}\prod_{j=1}^d(a_j,b_j-1/m]\). If \(x\in Q\), then
\(s=\min_j(b_j-x_j)>0\); choose \(m\) with \(1/m\leq s\), so that
\(x_j\leq b_j-1/m\) for every \(j\) and \(x\) lies in the \(m\)th
member of the union. Conversely, \(a_j<x_j\leq b_j-1/m<b_j\) gives
\(x\in Q\). The union is countable and its members lie in
\(\mathcal H_d\), so \(Q\in\sigma(\mathcal H_d)\). By
\cref{thm:p15-borel-generators} the rational open rectangles generate
\(\mathcal B(\mathbb R^d)\), so minimality gives
\(\mathcal B(\mathbb R^d)\subseteq\sigma(\mathcal H_d)\).

(3) By~(1), \(\mathcal H_d\) is a pi-system on \(\mathbb R^d\), and
by~(2) it generates \(\mathcal B(\mathbb R^d)\).
\Cref{thm:p15-measure-uniqueness}, applied with \(S=\mathbb R^d\) and
\(\mathcal C=\mathcal H_d\), shows that two probability measures that
agree on \(\mathcal H_d\) agree on \(\mathcal B(\mathbb R^d)\).
\end{proof}
\end{annexitem}

\begin{annexitem}{annex:p15-w02c-p15-014-point-mass}{thm:p15-w02c-p15-014-point-mass-law}{Point-mass laws and discrete laws}
\begin{proof}
Singletons are Borel: the complement of \(\{x\}\) is open, since a
point \(y\neq x\) has the ball of radius \(|y-x|\) about it inside the
complement, and sigma-algebras are closed under complements. A countable
set is a countable union of singletons and is Borel as well.

The law \(\mathbb P_X\) is a probability measure on
\(\mathcal B(\mathbb R)\) (\cref{thm:p15-pushforward-laws}).

(1) \emph{Measure.} \(\delta_c\) takes values in \(\{0,1\}\), and
\(\delta_c(\varnothing)=0\), \(\delta_c(\mathbb R)=1\). Let
\(B_1,B_2,\ldots\) be pairwise disjoint Borel sets with union \(B\). If
\(c\notin B\), then \(c\) lies in no \(B_n\), and both
\(\delta_c(B)\) and \(\sum_n\delta_c(B_n)\) are zero. If \(c\in B\),
then \(c\) lies in exactly one \(B_n\), by disjointness, so the series
has exactly one term equal to one and the others zero, and both sides
equal one. So \(\delta_c\) is countably additive, and it is a
probability measure (\cref{def:p15-measure}).

\emph{Distribution function.} \(F_c(x)=\delta_c((-\infty,x])\), which
is one exactly when \(c\leq x\) and zero when \(x<c\). At \(c=0\) this
is \(F_0(x)=1\) for \(x\geq0\) and \(F_0(x)=0\) for \(x<0\), Durrett's
function in Example~1.2.8.

\emph{Discrete.} \(\{c\}\) is countable, and
\(\delta_c(\mathbb R\setminus\{c\})=0\) because \(c\notin\mathbb
R\setminus\{c\}\).

\emph{Uniqueness.} Let \(\mu\) be a probability measure on
\(\mathcal B(\mathbb R)\) with \(\mu((-\infty,x])=F_c(x)\) for every
real \(x\). The rays \((-\infty,a]\), \(a\in\mathbb R\), form a
pi-system, because
\((-\infty,a]\cap(-\infty,b]=(-\infty,\min(a,b)]\), and they generate
\(\mathcal B(\mathbb R)\) (\cref{thm:p15-borel-generators}). The
measures \(\mu\) and \(\delta_c\) agree on every ray, so they agree on
\(\mathcal B(\mathbb R)\) by \cref{thm:p15-measure-uniqueness}.

(2) If \(\mathbb P_X=\delta_c\), then
\(p_X(c)=\mathbb P_X(\{c\})=\delta_c(\{c\})=1\). The statements
\(p_X(c)=1\) and ``\(X=c\) almost surely'' say the same thing, since
\(p_X(c)=\mathbb P(X=c)\) and \(\{X=c\}=X^{-1}(\{c\})\) is an event
(\cref{def:p15-measure}). Finally suppose \(p_X(c)=1\) and let \(B\) be
Borel. If \(c\in B\), then \(\{X=c\}\subseteq\{X\in B\}\), so
\(1\geq\mathbb P_X(B)\geq p_X(c)=1\). If \(c\notin B\), then
\(\{X\in B\}\subseteq\{X\neq c\}\), whose probability is
\(1-p_X(c)=0\) by \cref{thm:p15-measure-laws}~(1). In both cases \(\mathbb P_X(B)=\delta_c(B)\), so
\(\mathbb P_X=\delta_c\).

(3) The points of positive point mass are the discontinuities of
\(F_X\), and this set is at most countable
(\cref{thm:p15-cdf-laws}); so \(D\) is at most countable and Borel.

\emph{A sum formula.} Let \(C\) be an at most countable Borel set. List
its members as \(x_1,x_2,\ldots\), padding a finite list with empty
sets in place of further singletons. The singletons \(\{x_n\}\) are
pairwise disjoint with union \(C\), so countable additivity and
\cref{thm:p15-countable-atomic} (grouping of nonnegative terms, which
makes the sum independent of the listing) give
\[
 \mathbb P_X(C)=\sum_{x\in C}p_X(x).
\]

Suppose \(\mathbb P_X\) is discrete, with countable \(S\) and
\(\mathbb P_X(S^c)=0\). A point \(x\notin S\) has
\(p_X(x)\leq\mathbb P_X(S^c)=0\) by monotonicity
(\cref{thm:p15-measure-laws}), so \(D\subseteq S\). For Borel \(B\),
write \(B\) as the disjoint union of \(B\cap S\) and \(B\setminus S\).
The second part has \(\mathbb P_X(B\setminus S)\leq\mathbb P_X(S^c)=0\).
The sum formula with \(C=B\cap S\) gives
\(\mathbb P_X(B\cap S)=\sum_{x\in B\cap S}p_X(x)\), and the terms with
\(x\notin D\) are zero, so by grouping (\cref{thm:p15-countable-atomic})
the sum equals \(\sum_{x\in D\cap B}p_X(x)\). Hence
\(\mathbb P_X(B)=\sum_{x\in D\cap B}p_X(x)\). Taking \(B=\mathbb R\)
gives \(\sum_{x\in D}p_X(x)=1\).

Conversely, if \(\sum_{x\in D}p_X(x)=1\), the sum formula with \(C=D\)
gives \(\mathbb P_X(D)=1\), so \(\mathbb P_X(D^c)=0\) and \(\mathbb P_X\)
is discrete with the countable set \(D\).

For the last equality, \(\delta_x(B)\) is one for \(x\in B\) and zero
otherwise, so the terms of \(\sum_{x\in D}p_X(x)\delta_x(B)\) are
\(p_X(x)\) for \(x\in D\cap B\) and zero for the other \(x\in D\);
by grouping, the sum equals \(\sum_{x\in D\cap B}p_X(x)\).
\end{proof}
\end{annexitem}

\section{Uniform integrability: proofs}
\label{sec:annex-p15-w02c-p15-014-ui-source}

For a nonnegative random variable \(U\) and \(M\geq0\), the events
\(\{U\leq M\}\) and \(\{U>M\}\) partition \(\Omega\), so
\[
 EU=E(U;U\leq M)+E(U;U>M)
\]
in \([0,\infty]\) (\cref{thm:p15-nonnegative-integral-laws}~(4)), and
\(E(U;U\leq M)\leq M\,P(U\leq M)\leq M\) by monotonicity
(\cref{thm:p15-nonnegative-integral-laws}~(2)), since
\(U\mathbf1_{\{U\leq M\}}\leq M\mathbf1_{\{U\leq M\}}\). Both proofs below
apply this split to \(U=|X_i|\).

\begin{annexitem}{annex:p15-w02c-p15-014-ui-arbitrary}{thm:p15-w02c-p15-014-ui-arbitrary}{A level split}
\begin{proof}
(1) For each \(i\), the split gives
\(E|X_i|\leq M+E(|X_i|;|X_i|>M)\leq M+1\). Taking the supremum over
\(i\) gives \(\sup_iE|X_i|\leq M+1\). If \(I\) is empty, the supremum is
zero by the convention of \cref{def:p15-uniform-integrability}.

(2) Let \(M\geq0\). The split gives
\(\infty=E|X_j|=E(|X_j|;|X_j|\leq M)+E(|X_j|;|X_j|>M)\), and the first
term is at most \(M<\infty\), so the second term is \(\infty\). Hence
\(\sup_iE(|X_i|;|X_i|>M)=\infty\) for every \(M\geq0\), and the limit as
\(M\to\infty\) is not zero.

(3) Suppose the collection is uniformly integrable in Durrett's sense.
The limit being zero, there is \(M\geq0\) with
\(\sup_iE(|X_i|;|X_i|>M)\leq1\), and~(1) gives \(E|X_i|\leq M+1<\infty\)
for every \(i\). So every \(X_i\) is integrable, and the displayed
condition of \cref{def:p15-uniform-integrability}, which is the same
condition, holds. Conversely, if every \(X_i\) is integrable and the
family is uniformly integrable in the sense of
\cref{def:p15-uniform-integrability}, that definition's condition is
Durrett's condition.
\end{proof}
\end{annexitem}

\begin{annexitem}{annex:p15-w02c-p15-014-ui-moment}{thm:p15-w02c-p15-014-ui-moment-source}{Integrability and the tail bound from one moment bound}
\begin{proof}
(1) Since \(\varphi(x)/x\to\infty\), there is \(M_0>0\) with
\(\varphi(x)/x\geq1\) for every \(x\geq M_0\); for such \(x\),
\(\varphi(x)\geq x>0\). For \(M\geq M_0\) every \(x\geq M\) has
\(\varphi(x)>0\) and \(0<x/\varphi(x)\leq1\), so the set whose supremum
defines \(\varepsilon_M\) is a nonempty set of numbers in \((0,1]\) and
\(\varepsilon_M\in[0,1]\). If \(M_0\leq M\leq M'\), the set for \(M'\) is
contained in the set for \(M\), so \(\varepsilon_{M'}\leq\varepsilon_M\).
Given \(\eta>0\), there is \(M_\eta\geq M_0\) with
\(\varphi(x)/x\geq1/\eta\) for all \(x\geq M_\eta\); then
\(x/\varphi(x)\leq\eta\) for those \(x\), and \(\varepsilon_M\leq\eta\)
for every \(M\geq M_\eta\). Hence \(\varepsilon_M\to0\).

(2) Let \(M\geq M_0\) and \(i\in I\). At an outcome with
\(x=|X_i(\omega)|>M\), the number \(\varphi(x)\) is positive and
\[
 x=\frac{x}{\varphi(x)}\,\varphi(x)\leq\varepsilon_M\,\varphi(x),
\]
since \(x\geq M\) puts \(x/\varphi(x)\) in the set defining
\(\varepsilon_M\). At every other outcome the indicator below is zero.
So, pointwise on \(\Omega\),
\[
 |X_i|\mathbf1_{\{|X_i|>M\}}
 \leq\varepsilon_M\,\varphi(|X_i|)\mathbf1_{\{|X_i|>M\}}
 \leq\varepsilon_M\,\varphi(|X_i|).
\]
All three functions are nonnegative random variables, because
\(\varphi(|X_i|)\) is one by hypothesis and \(\{|X_i|>M\}\) is an event.
Monotonicity and homogeneity of the nonnegative integral
(\cref{thm:p15-nonnegative-integral-laws}) give
\(E(|X_i|;|X_i|>M)\leq\varepsilon_M E\varphi(|X_i|)\leq C\varepsilon_M\).

(3) Apply the split with \(M=M_0\) and then~(2) and~(1):
\[
 E|X_i|\leq M_0+E(|X_i|;|X_i|>M_0)\leq M_0+C\varepsilon_{M_0}\leq M_0+C<\infty .
\]
So every \(X_i\) is integrable. By~(2),
\(0\leq\sup_iE(|X_i|;|X_i|>M)\leq C\varepsilon_M\) for \(M\geq M_0\), and
the right side tends to zero by~(1). This is Durrett's condition, and
since every member is integrable it is also the condition of
\cref{def:p15-uniform-integrability}
(\cref{thm:p15-w02c-p15-014-ui-arbitrary}~(3)).
\end{proof}
\end{annexitem}

\section{Conditional limits: proofs}
\label{sec:annex-p15-w02c-p15-014-conditional}

\begin{annexitem}{annex:p15-w02c-p15-014-posterior-mean}{thm:p15-w02c-p15-014-posterior-mean}{The posterior mean is the conditional mean}
\begin{proof}
\emph{Integrability.} For the nonnegative measurable function
\(f(x)=|x|\), \cref{thm:p15-kernel-expectation} gives
\(\int|x|\,K_n(\cdot,dx)=E(|\theta|\mid\mathcal F_n)\) almost surely, in
the sense of the nonnegative extension of \cref{thm:p15-conditional-limits}.
Since \(E|\theta|<\infty\), the integrable definition applies, and its
version satisfies
\(E\bigl(E(|\theta|\mid\mathcal F_n)\bigr)=E|\theta|<\infty\) (the
all-event identity with the event \(\Omega\)). A nonnegative variable \(V\) with finite expectation is finite almost
surely: for every \(M>0\), monotonicity of the integral gives
\(EV\geq E(M\mathbf1_{\{V=\infty\}})=M\,P(V=\infty)\), so
\(P(V=\infty)>0\) would force \(EV=\infty\). Hence \(\int|x|\,K_n(\cdot,dx)<\infty\) almost surely.

\emph{Version.} Apply \cref{thm:p15-kernel-expectation} to the real
function \(f(x)=x\), with \(E|f(\theta)|=E|\theta|<\infty\). On the
full-measure set where the absolute kernel integral is finite, the signed
integral \(\int x\,K_n(\cdot,dx)\) equals \(E(\theta\mid\mathcal F_n)\),
and extending it by 0 on the measurable exceptional set gives an
integrable \(\mathcal F_n\)-measurable version. That extension is
\(m_n\). So \(m_n=E(\theta\mid\mathcal F_n)\) almost surely.

\emph{Convergence.} The hypotheses of
\cref{thm:p15-posterior-mean-consistency} hold with \(\xi_i=Z_i\): the
\(Z_i\) are i.i.d.\ and integrable with mean \(\mu=EZ_1\), \(\theta\) is
integrable and independent of the sequence, and \(Y_i=\theta+Z_i\). It
gives a version of \(E(\theta\mid\mathcal F_n)\) converging to
\(\theta\) almost surely and in \(L^1\). Two versions of
\(E(\theta\mid\mathcal F_n)\) agree almost surely, so for each \(n\),
\(m_n\) equals that version outside a null event \(N_n\). Outside the
null event \(\bigcup_nN_n\) (\cref{thm:p15-measure-laws}~(2)) the two
sequences coincide for all \(n\), so \(m_n\to\theta\) almost surely, and
\(E|m_n-\theta|\) equals the corresponding expectation for the version,
which tends to 0.
\end{proof}
\end{annexitem}
